%% Define your own pgfplot and tikz default styles
% This file defines TikZ styles and options for exams! It is also used be the ant_exam class!
\usepgfplotslibrary{groupplots}
\usetikzlibrary{arrows,arrows.meta,automata,backgrounds,calc,shapes.misc,shapes.symbols,shapes.geometric,decorations.markings,decorations.pathmorphing,decorations.shapes,matrix,positioning,shadows,svg.path,plotmarks,spy,pgfplots.groupplots,fit,patterns}
\pgfplotsset{compat=newest}

\tikzstyle{line} = [draw, -latex']
\tikzset{
	conn/.style={circle, fill=black, inner sep=.1pt, draw, minimum height = 1mm},
	box/.style={rectangle, fill=white, inner sep=7pt, very thick, draw, drop shadow={fill=white, xshift=0mm, yshift=-0mm}},
	kreis/.style={circle, fill=white, inner sep=7pt, very thick, draw,drop shadow={fill=white, xshift=0mm, yshift=-0mm}},
	se/.style={-stealth, thick},
	plus/.style={oplus,circle, fill=white, inner sep=.5pt, thick, draw, minimum height=3.5mm},
	times/.style={otimes,circle, inner sep=.5pt, thick, draw, minimum height=3.5mm},
	oplus/.style={path picture={\draw[black](path picture bounding box.south) -- (path picture bounding box.north) (path picture bounding box.west) -- (path picture bounding box.east);}},
	otimes/.style={path picture={\draw[black](path picture bounding box.south west) -- (path picture bounding box.north east) (path picture bounding box.south east) -- (path picture bounding box.north west);}},
	antenna/.style={isosceles triangle,fill=black,shape border rotate=-90, inner sep=2pt},
  	invisible/.style={opacity=0},
  	visible on/.style={alt=#1{}{invisible}},
  	alt/.code args={<#1>#2#3}{%
    	\alt<#1>{\pgfkeysalso{#2}}{\pgfkeysalso{#3}} % \pgfkeysalso doesn't change the path
  	},
	}

\tikzstyle{pluscircle} = [draw, circle, path picture={\draw (path picture bounding box.south) -- (path picture bounding box.north) (path picture bounding box.west) -- (path picture bounding box.east);}]

\tikzstyle{block} = [draw, rectangle, minimum height=2em, minimum width=4em, align=center, inner sep=2pt]
\tikzstyle{backgroundblock} = [minimum height=2em, minimum width=6em, align=center,fill=gray!30, inner sep=4pt]

\tikzstyle{input} = [coordinate]
\tikzstyle{output} = [coordinate]


% Define the ORCID LOGO
\usepackage{scalerel}
\usetikzlibrary{svg.path}
\definecolor{orcidlogocol}{HTML}{A6CE39}
\tikzset{
	orcidlogo/.pic={
		\fill[orcidlogocol] svg{M256,128c0,70.7-57.3,128-128,128C57.3,256,0,198.7,0,128C0,57.3,57.3,0,128,0C198.7,0,256,57.3,256,128z};
		\fill[white] svg{M86.3,186.2H70.9V79.1h15.4v48.4V186.2z}
		svg{M108.9,79.1h41.6c39.6,0,57,28.3,57,53.6c0,27.5-21.5,53.6-56.8,53.6h-41.8V79.1z M124.3,172.4h24.5c34.9,0,42.9-26.5,42.9-39.7c0-21.5-13.7-39.7-43.7-39.7h-23.7V172.4z}
		svg{M88.7,56.8c0,5.5-4.5,10.1-10.1,10.1c-5.6,0-10.1-4.6-10.1-10.1c0-5.6,4.5-10.1,10.1-10.1C84.2,46.7,88.7,51.3,88.7,56.8z};
	}
}

\newcommand\orcidicon[1]{\href{https://orcid.org/#1}{\mbox{\scalerel*{
				\begin{tikzpicture}[yscale=-1,transform shape]
				\pic{orcidlogo};
				\end{tikzpicture}
			}{U}}}} % original: |

% Official colors of the University of Bremen
\definecolor{UBBlueDark}{RGB}{28,53,107}
\definecolor{UBBlueMed}{RGB}{13,104,176}
\definecolor{UBBlueVeryLight}{RGB}{192,209,226} % this one's not official
\definecolor{UBBlueLight}{RGB}{114,179,223}
\definecolor{UBOrangeDark}{RGB}{247,167,3}
\definecolor{UBOrangeLight}{RGB}{255,232,177}
\definecolor{UBRedMed}{RGB}{213,17,48}
\definecolor{UBRedDark}{RGB}{135,39,70}

%%%%
% Carsten Style
%%%%
% A custom phantom function to manipulate the position of exponents in
% tikz plots
\def\myphantom#1#2{
  \setbox0 = \hbox{$#1$}
  \hbox to #2\wd0{\hfill}
}

% Insert tikz package for pictures and define library packages
\usetikzlibrary{arrows,calc,shapes.geometric,decorations.markings,decorations.pathmorphing,decorations.shapes,matrix,positioning,shadows,plotmarks,spy,pgfplots.groupplots,fit,patterns}
\pgfplotsset{compat=newest}

% Custom definitions to tweak standard plot formatting, plot groups
% and so on
\pgfplotsset{
  plot coordinates/math parser=false,
  log base 10 number format code/.code={$\displaystyle\ifthenelse{\equal{#1}{0.0}}{10^{\scriptscriptstyle\myphantom{\scriptscriptstyle-}{.3}\pgfmathprintnumber{#1}}\myphantom{\scriptscriptstyle-}{.7}}{10^{\scriptscriptstyle\pgfmathprintnumber{#1}}}$},
  every axis y label/.style={at={(ticklabel cs:0.5)},rotate=90,left=1mm,anchor=center},
  every axis x label/.style={at={(ticklabel cs:0.5)},below=1mm,anchor=center},
  every axis/.append style={
      footnotesize,
      width = 0.5\textwidth,
      legend style={font=\scriptsize,legend cell align = left,legend pos =north west},
      tick style={semithick},
  },
  every axis plot/.append style={thick},
  group style={
      columns= 2,
      horizontal sep=1.15cm,
      every plot/.style= {width=0.53\textwidth},
  },
  filter discard warning=false,
}

% Custom styles for block diagrams, especially times, plus,
% modulations in a node, etc.
\tikzset{
	common/.style={font=\small},
	box/.style={rectangle, fill=white, inner sep=7pt, very thick, draw,drop shadow={fill=black, xshift=1mm, yshift=-1mm}},
	vbox/.style={box,rotate=90,text centered},
	conn/.style={circle, fill=black, inner sep=.1pt, draw, minimum height = 1mm},
	oplus/.style={
          path picture={
            \draw[black](path picture bounding box.south) -- (path picture bounding box.north)
            (path picture bounding box.west) -- (path picture bounding box.east);
          }
        },
	otimes/.style={
          path picture={
            \draw[black](path picture bounding box.south west) -- (path picture bounding box.north east)
            (path picture bounding box.south east) -- (path picture
            bounding box.north west);
          }
        },
	4qam/.style={box,minimum height=8mm,minimum width=8mm,
          path picture={
            \draw[black,thin,-stealth](0,-3.1mm) -- (0,3.5mm);
            \draw[black,thin,-stealth] (-3.1mm,0) -- (3.5mm,0);
            \foreach \x in {-1.414213562373095mm,1.414213562373095mm} {
              \foreach \y in {-1.414213562373095mm,1.414213562373095mm} {
                \draw[fill=black,thin] (\x,\y) circle
                [radius=.1mm];
              }
            }
          }
        },
	16qam/.style={box,minimum height=8mm,minimum width=8mm,
          path picture={
            \draw[black,thin,-stealth](0,-3.1mm) -- (0,3.5mm);
            \draw[black,thin,-stealth] (-3.1mm,0) -- (3.5mm,0);
            \foreach \x in {-0.632455532033676mm,0.632455532033676mm,-1.897366596101028mm,1.897366596101028mm}{
              \foreach \y in {-0.632455532033676mm,0.632455532033676mm,-1.897366596101028mm,1.897366596101028mm}
              {
                \draw[fill=black,thin] (\x,\y) circle
                [radius=.1mm];
              }
            }
          }
        },
	64qam/.style={box,minimum height=8mm,minimum width=8mm,
          path picture={
            \draw[black,thin,-stealth](0,-3.1mm) -- (0,3.5mm);
            \draw[black,thin,-stealth] (-3.1mm,0) -- (3.5mm,0);
            \foreach \x in {-0.308606699924184mm,0.308606699924184mm,-0.925820099772551mm,0.925820099772551mm,-1.543033499620919mm,1.543033499620919mm,-2.160246899469286mm,2.160246899469286mm}{
              \foreach \y in {-0.308606699924184mm,0.308606699924184mm,-0.925820099772551mm,0.925820099772551mm,-1.543033499620919mm,1.543033499620919mm,-2.160246899469286mm,2.160246899469286mm}
              {
                \draw[fill=black,thin] (\x,\y) circle [radius=.1mm];
              }
            }
          }
        },
	plus/.style={oplus,circle, fill=white, inner sep=.5pt, thick, draw, minimum height=3.5mm},
	times/.style={otimes,circle, inner sep=.5pt, thick, draw, minimum height=3.5mm},
	antenna/.style={isosceles triangle,fill=black, shape border rotate=-90, inner sep=2pt},
	se/.style={-stealth, thick},
	pa/.style={-stealth,thick,double distance between line centers=2pt,
          decoration={markings,mark=at position .4 with {
              \draw (-2pt,-2pt) -- (2pt,2pt);\draw (0pt,-2pt) -- (4pt,2pt);
            }
          }
        },
	dots/.style={dash pattern=on 1pt off 4pt,very thick},
	wave/.style={-stealth,decorate,
          decoration={coil,aspect=0,amplitude=.2mm,segment length=1mm,pre=lineto,pre length=1mm,post=lineto, post length=1mm}
        },
	exit_inner/.style={color=black,solid},
	exit_combine/.style={color=black,dash pattern=on 1.5pt off 3pt on 1.5pt off 3pt, line width=1.5pt},
	exit_outer_bg/.style={color=white!50!black!50,dashed},
	exit_outer/.style={color=black,dashed},
}
% Declaration of layers in pgf pictures, e.g., to draw something
% behind a picture (also useful for beamer)
\pgfdeclarelayer{background}
\pgfsetlayers{background,main}
